% concatenated from VietorisThickeningsAndComplexesOfManifoldsAreHomotopyEquivalent.tex
\documentclass[11pt, reqno, english]{amsart}


\usepackage[utf8]{inputenc}
\usepackage[T1]{fontenc}
\usepackage{amsmath,amsthm}
\usepackage{amsfonts,amssymb}
\usepackage{url}
\usepackage{mathtools}  
\usepackage[colorlinks=true,urlcolor=blue,linkcolor=blue,citecolor=blue]{hyperref}
\usepackage{enumerate, paralist}
\usepackage{tikz,tikz-cd}
\usetikzlibrary{calc,arrows}
\usetikzlibrary{math} % Using variables inside tikz
\usetikzlibrary{arrows,decorations,backgrounds,calc,math}
\tikzstyle{map}=[->,semithick]
\tikzstyle{arc}=[bend left,->,semithick]
\tikzstyle{rinclusion}=[right hook->,semithick]
\tikzstyle{linclusion}=[left hook->,semithick]

\usepackage[left=1in,right=1in,top=1in,bottom=1in]{geometry}
\usepackage{xcolor}
\usepackage{hyperref}
\usepackage{float}
\usepackage{subfigure}
\usepackage{color}

\newtheoremstyle{myremark} % name
{7pt}                    % Space above
{7pt}                    % Space below
{}  	                 % Body font
{}                           % Indent amount
{\bf}       	         % Theorem head fuont
{.}                          % Punctuation after theorem head
{.5em}                       % Space after theorem head
{}  % Theorem head spec (can be left empty, meaning ‘normal’)

\theoremstyle{plain}
\newtheorem{lemma}{Lemma}[section]
\newtheorem{theorem}[lemma]{Theorem}
\newtheorem{corollary}[lemma]{Corollary}
\newtheorem{proposition}[lemma]{Proposition}
\newtheorem{conjecture}[lemma]{Conjecture}
\newtheorem{claim}[lemma]{Claim}
\theoremstyle{definition}
\newtheorem{definition}[lemma]{Definition}
\newtheorem*{definition-cnk}{Definition~\ref{def:cnk}}
\newtheorem{question}{Question}
\newtheorem*{question-motivating}{Motivating Question}
\newtheorem{fact}[lemma]{Fact}
\makeatletter
\newenvironment{subtheorem}[1]{%
  \def\subtheoremcounter{#1}%
  \refstepcounter{#1}%
  \protected@edef\theparentnumber{\csname the#1\endcsname}%
  \setcounter{parentnumber}{\value{#1}}%
  \setcounter{#1}{0}%
  \expandafter\def\csname the#1\endcsname{\theparentnumber\Alph{#1}}%
  \ignorespaces
}{%
  \setcounter{\subtheoremcounter}{\value{parentnumber}}%
  \ignorespacesafterend
}
\makeatother
\newcounter{parentnumber}

\theoremstyle{myremark}
\newtheorem{remark}[lemma]{Remark}
\newtheorem{example}[lemma]{Example}

\newcommand{\C}{\mathbb{C}}
\newcommand{\N}{\mathbb{N}}
\newcommand{\Q}{\mathbb{Q}}
\newcommand{\R}{\mathbb{R}}
\newcommand{\Z}{\mathbb{Z}}
\newcommand{\F}{\mathbb{F}}
\newcommand{\RP}{\ensuremath{\mathbb{R}\mathrm{P}}}
\newcommand{\CP}{\ensuremath{\mathbb{C}\mathrm{P}}}
\newcommand{\cA}{\mathcal{A}}
\newcommand{\cB}{\mathcal{B}}
\newcommand{\cF}{\mathcal{F}}
\newcommand{\cM}{\mathcal{M}}
\newcommand{\cN}{\mathcal{N}}
\newcommand{\cK}{\mathcal{K}}
\newcommand{\cP}{\mathcal{P}}
\newcommand{\cU}{\mathcal{U}}
\newcommand{\cV}{\mathcal{V}}
\newcommand{\cW}{\mathcal{W}}
\newcommand{\cY}{\mathcal{Y}}
\newcommand{\sP}{\mathscr{P}}
\newcommand{\image}{\mathrm{im}}
\newcommand{\diam}{\mathrm{diam}}
\newcommand{\cost}{\mathrm{cost}}
\newcommand{\dis}{\mathrm{dis}}
\newcommand{\dist}{\mathrm{dist}}
\newcommand{\cl}{\mathrm{cl}}
\newcommand{\conn}{\mathrm{conn}}
\newcommand{\conv}{\mathrm{conv}}
\newcommand{\cone}{\mathrm{cone}}
\newcommand{\st}{\mathrm{st}}
\newcommand{\sm}{\mathrm{SM}}
\newcommand{\m}{\mathrm{M}}
\newcommand{\SM}{\mathbb{SM}}
\newcommand{\M}{\mathbb{M}}
\newcommand{\sym}{\mathrm{Sym}}
\newcommand{\lip}{\mathrm{Lip}}
\newcommand{\sd}{\mathrm{sd}}

\newcommand{\Conf}{\mathrm{Conf}}

\newcommand{\id}{\mathrm{id}}
\newcommand{\oo}{\mathrm{O}}
\newcommand{\so}{\mathrm{SO}}
\newcommand{\gh}{\mathrm{GH}}
\newcommand{\h}{\mathrm{H}}
\newcommand{\ph}{\mathrm{PH}}
\newcommand{\sign}{\mathrm{sign}}
\newcommand{\im}{\mathrm{im}}
\newcommand\abs[1]{\left|#1\right|}
\newcommand{\supp}{\mathrm{supp}}
\newcommand{\ut}{\underline{t}}
\newcommand{\ind}{\mathrm{ind}}
\newcommand{\coind}{\mathrm{coind}}
\newcommand{\codis}{\mathrm{codis}}
\newcommand{\cov}{\mathrm{cov}}
\newcommand{\numCover}{\mathrm{numCover}}
\newcommand{\pack}{\mathrm{pack}}
\newcommand{\vol}{\mathrm{vol}}
\newcommand{\eps}{\varepsilon}
\newcommand{\ord}{\mathrm{ord}}
\newcommand{\card}{\mathrm{card}}
\newcommand{\comp}{\mathrm{Comp}}

\newcommand{\note}[1]{\textcolor{blue}{({#1})}}

\newcommand{\vr}[2]{\mathrm{VR}(#1;#2)}
\newcommand{\viet}[1]{\mathrm{V}(#1)}
\newcommand{\vietm}[1]{\mathrm{V}^\mathrm{m}(#1)}
\newcommand{\vrless}[2]{\mathrm{VR}_<(#1;#2)}
\newcommand{\vrleq}[2]{\mathrm{VR}_\le(#1;#2)}
\newcommand{\vrm}[2]{\mathrm{VR}^\mathrm{m}(#1;#2)}
\newcommand{\vrmless}[2]{\mathrm{VR}^\mathrm{m}_<(#1;#2)}
\newcommand{\vrmleq}[2]{\mathrm{VR}^\mathrm{m}_\le(#1;#2)}
\newcommand{\cech}[2]{\mathrm{\check{C}}(#1;#2)}
\newcommand{\acech}[3]{\mathrm{\check{C}}(#1,#2;#3)}
\newcommand{\cechm}[2]{\mathrm{\check{C}}^\mathrm{m}(#1;#2)}
\newcommand{\acechm}[3]{\mathrm{\check{C}}^\mathrm{m}(#1,#2;#3)}

\DeclareMathOperator{\sk}{\mathrm{sk}}
\DeclareMathOperator{\relint}{\mathrm{relint}}

\newcommand{\cpl}{\ensuremath{\mathrm{Cpl}}}

\linespread{1.2}

\usepackage{accents}
\usepackage[normalem]{ulem}
\usepackage{wasysym}
% The following three commands are used for the new font,
% uncomment to use the usual font
\usepackage{charter}
\usepackage{euler}
\usepackage[T1]{fontenc}
%
\usepackage{thmtools}
\usepackage{thm-restate}
\newcommand{\og}{\overline{g}}
\newcommand{\oh}{\overline{h}}
\newcommand{\of}{\overline{f}}
\newcommand{\ophi}{\overline{\phi}}
\newcommand{\opsi}{\overline{\psi}}
\def\fr{\mathrm{FillRad}}


\begin{document}

\title{Vietoris thickenings and complexes of manifolds are homotopy equivalent}

\author{Henry Adams}
\author{Alexandre Karassev}
\author{\v{Z}iga Virk}

\begin{abstract}
We show that if $X$ is a finite-dimensional Polish metric space, then the natural bijection $\vr{X}{r}\to\vrm{X}{r}$ from the (open) Vietoris--Rips complex to the Vietoris--Rips metric thickening is a homotopy equivalence.
This occurs, for example, if $X$ is a Riemannian manifold.
The same is true for the map \v{C}$(X;r)\to$\v{C}$^\mathrm{m}(X;r)$
%$\cech{X}{r}\to\cechm{X}{r}$ \note{I don't know why this isn't compiling correctly} 
from the \v{C}ech complex to the \v{C}ech metric thickening, and more generally, for the natural bijection $\viet{\cW}\to\vietm{\cW}$ from the Vietoris complex to the Vietoris metric thickening of any uniformly bounded cover $\cW$ of a finite dimensional Polish metric space.
We also show that if $X$ is a compact metrizable space, then $\vietm{\cW}$ is strongly locally contractible.
\end{abstract}

\subjclass[2020]{55N31, 54F45, 54C55, 54C65}

\keywords{Vietoris--Rips complexes, \v{C}ech complexes, metric thickenings, dimension theory, absolute neighborhood retracts, Michael's selection theorem}

\maketitle

%\setcounter{tocdepth}{2}
%\tableofcontents



\section{Introduction}

Let $X$ be a metric space, and let $r>0$.
The \emph{Vietoris--Rips simplicial complex $\vr{X}{r}$} has $X$ as its vertex set, and a finite subset $\sigma\subseteq X$ as a simplex when its diameter is less than $r$.
Vietoris--Rips complexes were invented by %Leopold
Vietoris in a cohomology theory for metric spaces~\cite{Vietoris27,lefschetz1942algebraic}.
Independently, they were introduced by %Eliyahu
Rips in geometric group theory as a natural way to thicken a space.
Indeed, Rips used these complexes to show that torsion-free hyperbolic groups have Eilenberg--MacLane spaces that are finite simplicial complexes~\cite{bridson2011metric}.

More recently, Vietoris--Rips complexes and persistent homology have become commonly used tools in applied and computational topology, motivated by applications to
data analysis~\cite{CarlssonIshkhanovDeSilvaZomorodian2008,ghrist2008barcodes}.
%Since we do not know \emph{a priori} how to choose the scale $r$, the idea of persistent homology is to compute the homology of the Vietoris--Rips complex of a dataset $X$ over a large range of scale parameters $r$ and to trust those topological features which persist.
For $M$ a manifold and scale $r$ sufficiently small depending on the curvature of $M$, Hausmann~\cite{Hausmann1995} proves there is a homotopy equivalence $\vr{M}{r} \simeq M$.
Furthermore, for dataset $X$ sufficiently close to $M$ in the Gromov--Hausdorff distance, Latschev~\cite{Latschev2001} proves there is a homotopy equivalence $\vr{X}{r} \simeq M$.

If $X$ is not a discrete metric space, then the inclusion $X\hookrightarrow\vr{X}{r}$ is not continuous, since the vertex set of a simplicial complex is equipped with the discrete topology.
For example, if $M$ is a manifold of dimension at least one, then $M\hookrightarrow\vr{M}{r}$ is not continuous.
This observation helped motivate the definition of Vietoris--Rips metric thickenings in~\cite{AAF}.
For $X$ a metric space and $r>0$, the \emph{Vietoris--Rips metric thickening $\vrm{X}{r}$} is the space of finitely-supported probability measures on $X$ whose support has diameter less than $r$, equipped with an optimal transport metric.
The inclusion $X\hookrightarrow\vrm{X}{r}$, obtained by mapping a point $x\in X$ to the Dirac delta measure $\delta_x$, is now an isometric embedding.

What is the relationship between Vietoris--Rips simplicial complexes and metric thickenings?
For $X$ a finite metric space, $\vr{X}{r}$ and $\vrm{X}{r}$ are homeomorphic~\cite%[Proposition~6.2]
{AAF}.
For $X$ an arbitrary metric space, Gillespie showed that the natural bijection $\vr{X}{r}\to\vrm{X}{r}$ is a weak homotopy equivalence, i.e., induces isomorphisms on all homotopy groups~\cite{gillespie2024vietoris}, improving upon~\cite{HA-FF-ZV}.
The question remains if $\vr{X}{r}$ and $\vrm{X}{r}$ are homotopy equivalent.
Since $\vr{X}{r}$ is a simplicial complex (and hence a CW complex), this result would be implied by Whitehead's theorem if $\vrm{X}{r}$ were known to be an absolute neighborhood retract, %or an absolute neighborhood extensor, 
but to the best of our knowledge this question remains unknown for general metric spaces $X$.

We prove that if $X$ is a finite-dimensional Polish metric space and $r>0$, then
\begin{itemize}
\item $\vrm{X}{r}$ is an absolute neighborhood retract, and hence
\item the natural bijection $\vr{X}{r}\to\vrm{X}{r}$ from the (open) Vietoris--Rips complex to the Vietoris--Rips metric thickening is a homotopy equivalence.
\end{itemize}
This occurs, for example, if $X=M$ is a Riemannian manifold.

As a result, we get a version of Latschev's theorem for metric thickenings, resolving~\cite[Conjecture~6.9]{AAR} in the affirmative; see Corollary~\ref{cor:latschev}.

More generally, for $\cW$ a uniformly bounded cover of a Polish metric space, we show that the natural bijection $\viet{\cW}\to\vietm{\cW}$ from the Vietoris complex to the Vietoris metric thickening is a homotopy equivalence; see Theorem~\ref{thm:main}.
This also implies that the map $\cech{X}{r}\to\cechm{X}{r}$ from the (open) \v{C}ech complex to the \v{C}ech metric thickening is a homotopy equivalence.

Furthermore, if $X$ is a compact metrizable space, then we show that $\vietm{\cW}$ is strongly locally contractible, using Michael's selection theorem~\cite{michael1956continuous}.

We begin in Section~\ref{sec:preliminaries} with background and definitions.
In Section~\ref{sec:homotopy-equiv} we prove our main result, Theorem~\ref{thm:main}.
We prove that $\vietm{\cW}$ is strongly locally contractible for $X$ a compact metrizable space in Section~\ref{sec:strong-local-contractibility}.
In Section~\ref{sec:conclusion} we conclude with a list of open questions.




%\section{Related work}



\section{Preliminaries}
\label{sec:preliminaries}

We provide some background on metric spaces, Vietoris--Rips and \v{C}ech simplicial complexes, Vietoris simplicial complexes,  spaces of measures, metric thickenings, absolute neighborhood retracts, Whitehead's theorem, dimension, and local contractibility.


\subsection{Metric spaces}

Let $X=(X,d)$ be a metric space.
The \emph{diameter} of a subset $A\subseteq X$ is defined as $\diam(A)=\sup\{d(a,a')~|~a,a'\in A\}$.
The \emph{open ball} of radius $r>0$ centered at $x\in X$ is $B(x;r)=\{y\in X~|~d(y,x)<r\}$, which we sometimes denote as $B_X(x;r)$.


\subsection{Vietoris--Rips and \v{C}ech simplicial complexes}

Let $X$ be a metric space and let $r>0$.

The \emph{Vietoris--Rips simplicial complex $\vr{X}{r}$} has $X$ as its vertex set, and a finite subset $\sigma\subseteq X$ as a simplex when $\diam(\sigma)<r$.

The \emph{intrinsic \v{C}ech simplicial complex} $\cech{X}{r}$ has $X$ as its vertex set, and a finite subset $\sigma\subseteq X$ as a simplex when $\cap_{x\in \sigma}B(x;r)\neq\emptyset$.

If $Z\supseteq X$ is a metric space extending the metric on $X$, then the \emph{ambient \v{C}ech simplicial complex} $\acech{X}{Z}{r}$ has $X$ as its vertex set, and a finite subset $\sigma\subseteq X$ as a simplex when $\cap_{x\in \sigma}B_Z(x;r)\neq\emptyset$.

A common example of an ambient \v{C}ech complex is when $X\subseteq Z=\R^n$.
When $Z=X$, we note that the ambient \v{C}ech complex recovers the intrinsic \v{C}ech complex, namely $\acech{X}{X}{r}=\cech{X}{r}$.
For this reason, we can state our results with ambient \v{C}ech complexes while still handling the case of intrinsic \v{C}ech complexes.

We identify simplicial complexes with their geometric realizations.

\subsection{Vietoris simplicial complexes}
\label{ssec:vietoris}

Given a cover $\cW$ of a metric space, the \emph{Vietoris complex} $\viet{\cW}$ of $\cW$ has $X$ as its vertex set, and a finite subset $\sigma \subseteq X$ as a simplex when $\sigma$ is contained in an element of $\cW$.
By Dowker duality~\cite{dowker1952topology}, $\viet{\cW}$ is homotopy equivalent to the nerve of $\cW$.

Let $ r>0$.
When $\cW$ consists of all open subsets of diameter less than $r$, we have $\viet{\cW}=\vr{X}{r}$.
If instead $\cW=\{B_X(x;r)\}_{x\in X}$ consists of all open $r$-balls in $X$, then $\viet{\cW}=\cech{X}{r}$.
For $Z\supseteq X$ a metric space extending the metric on $X$, if $\cW=\{X\cap B_Z(z;r)\}_{z\in Z}$, then $\viet{\cW}=\acech{X}{Z}{r}$.


\subsection{Spaces of measures}

We follow~\cite{villani2008optimal,fedorchuk1991probability}.

%\note{Define a probability measure.}

For $X$ a topological space, we let $P(X)$ be the space of probability measures on the Borel sigma algebra of $X$, equipped with the weak topology.
The weak topology on $P(X)$ is induced by convergence against $C_b(X)$, the bounded continuous test functions.
In more detail, a sequence $\mu_1, \mu_2, \mu_3, \ldots$ in $P(X)$ is said to converge weakly to $\mu\in P(X)$ if for all bounded continuous functions $f\colon X\to \R$, we have $\lim_{n \to \infty} \int_X f(x)\,d\mu_n = \int_X f(x)\,d\mu$.

For $k\ge 0$, let $P_k(X)$ be the set of all measures in $P(X)$ whose support have cardinality at most $k$.
So, a measure $\mu \in P_k(X)$ can be written as $\mu=\sum_{i=0}^{k-1}\lambda_i \delta_{x_i}$ for some $x_i\in X$ and $\lambda_i \ge 0$ with $\sum_{i=0}^{k-1}\lambda_i=1$.
The \emph{support} of a measure $\mu = \sum_{i=0}^{k-1}\lambda_i \delta_{x_i} \in P_k(X)$ is $\supp(\mu)=\{x_i~|~\lambda_i>0\}$.

If $(X,d)$ is a metric space, then we may also define optimal transport metrics on $\cup_{k\ge \infty}P_k(X)$, the space of all finitely supported measures.
Given finitely supported measures $\mu_1$ and $\mu_2$, a \emph{coupling} between them is a probability measure $\mu$ on $X\times X$ with marginals $\mu_1$ and $\mu_2$.
This means $(\pi_1)_\sharp\mu=\mu_1$ and $(\pi_2)_\sharp \mu= \mu_2$, where $\pi_i\colon X\times X\to X$ is the projection map defined by $\pi_i(x_1,x_2)=x_i$ for $i=1,2$, and where $\sharp$ denotes the pushforward.
For $q\in[1,\infty)$ we may equip $\cup_{k\ge \infty}P_k(X)$ with the $q$-\emph{Wasserstein-Kantorovich-Rubinstein distance}, given by
\begin{equation}
d_{W^q}(\mu_1,\mu_2) = \inf_{\mu} \left(\int_{X\times X}(d(x,x'))^q\,d\mu\right)^{\frac{1}{q}},
\end{equation}
where $\mu$ varies over all couplings between $\mu_1$ and $\mu_2$.

%It is known that on a bounded metric space, for any $q \in [1, \infty)$, the $q$-Wasserstein metric generates the weak topology.
%For a summary of the result, see Appendix~\ref{app:Metrization_weak_topology}.
%On the other hand, the $\infty$-Wasserstein metric $\dWqq{\infty}$ generates a finer topology than the weak topology in general.

A \emph{Polish space} $X$ is a separable completely metrizable topological space, that is, a space homeomorphic to a complete metric space that has a countable dense subset.

%\note{For $q\in [1,\infty)$, define the $q$-Wasserstein-Kantarovoich-Rubenstein metric on $P(X)$.}

Let $k\ge 1$, and let $q\in [1,\infty)$.
If $X$ is a bounded Polish metric space, then the $q$-Wasserstein-Kantarovoich-Rubenstein metric on $P_k(X)$ agrees with the topology induced by weak convergence on $P_k(X)$ by~\cite[Corollary~6.13]{villani2008optimal}.
In turn, the topology induced by weak convergence coincides with the weak topology on $P(X)$ (see, e.g.,~\cite[Definitions~8.1.1 and 8.1.2]{bogachev2007measure} and~\cite{kallianpur1961topology}).
As a consequence, different choices of $q\in [1,\infty)$ all induce the same topology on $P_k(X)$, and we have the following observation.

\begin{proposition}
\label{proposition:weaktopology}
Let $X$ be a bounded Polish space.
Then the weak topology on $P_k(X)$ coincides with the topology induced by the $q$-Wasserstein-Kantarovoich-Rubenstein metric for any $q\in [1,\infty)$.
\end{proposition}
%\footnote{See Appendix A.1 with Memoli, Moy, Wang --- could instead add $X$ bounded assumption, but instead would be better to add $X$ Polish assumption.}

%\textcolor{orange}{Alex: I've made some edits related to the discussion below so I hope everything looks ok here now.}

%\note{This result seems pretty important and so we should probably state it explicitly as a citeable theorem.
%It means that throughout this paper, we can actually work with the weak topology, which is especially true for the last section, where we prove conractibility in the weak (i.e., not Wasserstein metric) topology.
%Also, does~\cite[Theorem~6.9]{villani2008optimal} require $X$ to be bounded?
%We should clarify the situation.
%Or we could prove directly that the Wasserstein and topology induced by sets $N(\mu, \mathcal{U}, \varepsilon )$ are equivalent.
%The fact that equivalence of convergences in two metrizable topologies imply homeomorphism follows from Theorem 21.3 of Munkres' Topology.}

%\textcolor{red}{Alex: It is sufficient for us for Wasserstein metric to agree with weak topology on $P_k$, since from it we just need the finite-dimensionality.
%However, I now have doubts that~\cite[Theorem~6.9]{villani2008optimal} would imply  that the weak topology on $P_k$ is the same as the one generated by the Wasserstein distance without the assumption of boundedness.
%Nevertheless, we still can prove that $P_k$ is finite-dimensional as follows: $P_k$ is the countable union of $P_k(n)$ where $P_k(n)$ is the set of all measures whose support consists of at most $k$ points and is contained in a closed ball $B(x_0,n)$, where $x_0$ is a fixed base point in $X$; the Wasserstein metric topology on $P_k(n)$ coincides with the weak topology because of Corollary~6.13  Villani's book.
%Therefore the dimension of $P_k(n)$ is bounded by some $M$ (that depends on dimension of $X$ and on $k$).
%Consequently, by the countable sum theorem, the dimension of $P_k$ (considered with the Wasserstein topology) is also bounded by $M$.
%For the last section, $X$ is compact so the required statement follows from Corollary~6.13.}


\subsection{Metric thickenings}

For $X$ a metric space and $r>0$, the \emph{Vietoris--Rips metric thickening} is the space
\[\vrm{X}{r}=\left\{\sum_{i=0}^{k-1}\lambda_i\delta_{x_i}\in P(X)~\mid~k\ge 1,\ \lambda_i\ge 0,\ \sum_i\lambda_i=1,\ \diam(\{x_0,\ldots,x_{k-1}\})<r\right\},\]
equipped with the optimal transport (or $1$-Wasserstein-Kantarovoich-Rubenstein) metric.

More generally, for an open cover $\cW$ of $X$, the \emph{metric thickening} of $\viet{\cW}$ is the space
\[
\vietm{\cW}=\left\{\sum_{i=0}^{k-1}\lambda_i\delta_{x_i}\in P(X)~\mid~k\ge 1,\ \lambda_i\ge 0,\ \sum_i\lambda_i=1,\ \exists W \in \cW: \{x_0,\ldots,x_{k-1}\} \subseteq W \right\},
\]
equipped with the optimal transport metric.
Choosing $\cW$ as in Section~\ref{ssec:vietoris} gives the \emph{\v{C}ech metric thickenings} $\cechm{X}{r}$ and $\acechm{X}{Z}{r}$.

The natural bijection $\vr{X}{r}\to \vrm{X}{r}$ is the continuous map defined by sending a point $\sum \lambda_i x_i$ in the geometric realization of $\vr{X}{r}$ to the probability measure $\sum \lambda_i\delta_{x_i}$ in $\vrm{X}{r}$.
The same is true for the natural bijection $\viet{\cW}\to \vietm{\cW}$.

In Corollary~\ref{cor:main} we show that for $X$ a finite-dimensional Polish metric space, the natural bijection $\vr{X}{r}\to\vrm{X}{r}$ is a homotopy equivalence.\footnote{
We remark this result is false if we had defined the Vietoris--Rips complex and metric thickening with the $\le$ convention instead of the $<$ convention, meaning a simplex is included in $\vr{X}{r}$ if its diameter is $\le r$ instead of $<r$.
Indeed, $\vrleq{X}{0}$ is $X$ with the discrete topology, whereas $\vrmleq{X}{0}$ is $X$ with its standard topology induced from its metric.
}


\subsection{Topological spaces}

For $Y$ a topological space and $y\in Y$, we let $\pi_n(Y,y)$ denote the $n$-th homotopy group of $Y$, with basepoint $y$.

A subset $A$ of a space $Y$ is a \emph{retract of $Y$} if there is a map $r\colon Y\to A$ with $r|_A=\id_A$.
This implies that $A$ is closed in $Y$.
A \emph{neighborhood retract of $Y$} is a closed set in $Y$
that is a retract of some neighborhood in $Y$.
A metrizable space $Y$ is called an
\emph{absolute neighborhood retract (ANR)} if $Y$ is a
neighborhood retract of any metrizable space that contains $Y$ as a closed subspace.
See~\cite{sakai2013geometric} for more information on ANRs.

By~\cite[Corollary~6.6.5]{sakai2013geometric} a space has the homotopy type of a simplicial complex if and only if it has the homotopy type of an ANR.


\subsection{Whitehead's theorem}

Let $Y$ and $Y'$ be topological spaces.
A map $f\colon Y\to Y'$ is a \emph{weak homotopy equivalence} if it induces
\begin{itemize}
\item a bijection from the path components of $Y$ to the path components of $Y'$, and
\item an isomorphism $\pi_n(Y,y)\to\pi_n(Y',y)$ for all $n\ge 1$ and $y\in Y$.
\end{itemize}

The classical version of Whitehead's theorem states that if $Y$ and $Y'$ are CW complexes, then every weak homotopy equivalence $f\colon Y\to Y'$ is a homotopy equivalence~\cite{Hatcher}.
%Simplicial complexes, including Vietoris--Rips simplicial complexes, are CW complexes.

The CW complex assumption can be relaxed, as follows.
By~\cite[Corollary~6.6.6]{sakai2013geometric}, if $Y$ and $Y'$ are ANRs, then every weak homotopy equivalence $f\colon Y\to Y'$ is a homotopy equivalence.

%We can combine Corollaries~6.6.5 and~6.6.6 of~\cite{sakai2013geometric} to get that if $K$ is a simplicial complex and $Y'$ is an ANR, then every weak homotopy equivalence $f\colon K\to Y'$ is a homotopy equivalence.
%Indeed, let $g\colon Y\to K$ be a homotopy equivalence from an ANR $Y$ to the simplicial complex $K$, which exists by~\cite[Corollaries~6.6.5]{sakai2013geometric}.
%Then $fg\colon Y\to Y'$ is a weak homotopy equivalence between ANRs, which is therefore a homotopy equivalence by~\cite[Corollaries~6.6.6]{sakai2013geometric}.
%Since $g$ and $fg$ are each homotopy equivalences, the two-out-of-three property~\note{\cite{}} then gives that $f$ is a homotopy equivalence, as desired.


\subsection{Dimension}

We refer the reader to Chapter~5 of Sakai's book~\cite{sakai2013geometric} for background on dimension.

Let $Y$ be a topological space, and let $\cU$ be an open cover of $Y$.
For any point $y\in Y$, let $\cU[y]=\{U\in\cU~|~y\in U\}$ be the subcollection of sets in $\cU$ containing $y$, and let $\card(\cU[y])$ be the cardinality of this set.
The \emph{order} of $\cU$ is $\ord(\cU)=\sup\{\card(\cU[y])~|~y\in Y\}$.
We define the \emph{(covering) dimension $\dim Y$ of $Y$} as follows.
Let $n$ be a nonnegative integer.
If every finite open cover of $Y$ has a finite open refinement $\cU$ with $\ord(\cU)\le n+1$, then $\dim Y\le n$.
(When $Y$ is paracompact, the word ``finite'' can be removed in both places in the prior sentence, still yielding an equivalent definition.)
We have $\dim Y=n$ if $\dim Y\le n$ and $\dim Y\nleq n-1$, in which case we say that the space $Y$ is \emph{$n$-dimensional}.
The space $Y$ is \emph{finite-dimensional} if $\dim Y\le n$ for some finite $n$, and otherwise $Y$ is \emph{infinite-dimensional}.

The space $Y$ is \emph{countable-dimensional} if it is a countable union of finite-dimensional normal subspaces.
%\note{Page 279 of Sakai writes:} It is said that $X$ is countable-dimensional (c.d.) if $X$ is a countable union of f.d.\ normal subspaces, where it should be noted that subspaces of normal spaces need not be normal (cf.\ Sect.\ 2.10).
Though a subspace of a normal space need not be normal, a subspace of a metric space is a metric space and hence is normal.
Therefore, a metric space $X$ is countable-dimensional if it is a countable union of finite-dimensional subspaces.

The countable sum theorem~\cite[Theorem~5.4.1]{sakai2013geometric} says that if the countable union $Y=\cup_{i\ge 1}F_i$ is normal, if each $F_i$ is closed in $Y$, and if $\dim F_i\le n$ for every $i\ge 1$, then $\dim Y\le n$.

Let $X$ be a metric space.
The subset theorem~\cite[Theorem~5.3.3]{sakai2013geometric} says that for every subset $A$ of a metrizable space $X$, we have $\dim A \le \dim X$.

The decomposition theorem~\cite[Theorem~5.4.5]{sakai2013geometric} states that a metric space $X$ satisfies $\dim X\le n$ if and only if $X$ is covered by $n+1$ many $0$-dimensional subsets $X_1$, \ldots, $X_{n+1}$.
It follows that a metric space $X$ is countable-dimensional if and only if it is a countable union of 0-dimensional subspaces.

% Recall from~\cite[Section~5.6]{sakai2013geometric} that $Y$ is \emph{countable-dimensional} if $Y=\cup_{i\ge 1} A_i$ for some countably many subsets $A_i\subseteq Y$ with $\dim A_i \le 0$.
%By~\cite[Theorem~5.4.5]{sakai2013geometric}, every $n$-dimensional metrizable space is the union of at most $n+1$ many $0$-dimensional subspaces.
%It follows that if $Y$ is the countable union of finite-dimensional subspaces, then $Y$ is countable-dimensional.

By~\cite[Theorem~6.10.1]{sakai2013geometric}, every countable-dimensional locally contractible metrizable space is an ANR.


\subsection{Local contractibility}
\label{ssec:local-contractibility}

A topological space $Y$ is \emph{locally contractible at $y \in Y$} if for every neighborhood $U$ of $y$, there is a neighborhood $y\in V\subseteq U$ such that the inclusion $V\hookrightarrow U$ is nullhomotopic in $U$.
The space $Y$ is \emph{locally contractible} if it is locally contractible at each $y\in Y$.

A topological space $Y$ is \emph{strongly locally contractible at $y \in Y$} if for every neighborhood $U$ of $y$, there is a contractible neighborhood $y\in V\subseteq U$.
In other words, if every point $y\in Y$ has a local base of contractible neighborhoods, then we say that $Y$ is \emph{strongly locally contractible}.

A neighborhood $U$ \emph{strongly} deformation retracts onto $y$ if the deformation retract fixes the point $y$ throughout the homotopy.
If every point $y\in Y$ has a local base of neighborhoods which strongly deformation retract onto $y$, then we say that $Y$ is \emph{strictly strongly locally contractible}.



\section{$\viet{\cW}\to\vietm{\cW}$ is a homotopy equivalence for $X$ finite-dimensional}
\label{sec:homotopy-equiv}

%Alex: If $\vrm{X}{r}$ is weakly infinite-dimensional (which should be implied if $X$ is finite-dimensional, such as a manifold), then strong local contractibility of $\vrm{X}{r}$ would give that $\vrm{X}{r}$ is an ANR.
%Hence $\vr{X}{r}\to\vrm{X}{r}$ would be a homotopy equivalence, (by Whitehead's theorem, since it is already known to be a weak homotopy equivalence).

%Theorem~0.42 on page 26 of (\footnote{\url{https://repovs.splet.arnes.si/files/2022/09/Continuous-Selections-of-Multivalued-Mappings-1998.pdf}}) works if $X$ is finite-dimensional, but we want a stronger version where $X$ is only assumed to be weakly infinite-dimensional.

%Metric space $X$ finite-dimensional, space of measures with at most $m$ points in support is finite-dimensional, so $\vrm{X}{r}$ is countably-dimensional.
%(Perhaps even strongly countably-dimensional, instead of weakly so.)

%Book: Katsuro Sakai, ``Geometric aspects of general topology'', page 392, theorem 6.10.1.

%Alex and Ziga are pointing out for $X$ finite-dimensional, (standard, not strong) local contractibility  along with Sakai's result is enougth to give that $\vr{X}{r}\to\vrm{X}{r}$ is a homotopy equivalence --- and we already know that $\vrm{X}{r}$ is locally contractible by our prior paper?

We prove that if $\cW$ is a uniformly bounded open cover of a finite-dimensional metric space $X$, then the Vietoris complex of $\cW$ and its metric thickening are homotopy equivalent.
A cover $\cW$ of $X$ is \emph{uniformly bounded} if there exists some $D>0$ such that $\diam(W) < D$ for all $W \in \cW$.

\begin{theorem}
\label{thm:main}
Let $\cW$ be a uniformly bounded open cover of a finite-dimensional Polish metric space $X$.
The natural bijection $\viet{\cW}\to\vietm{\cW}$ from the  Vietoris complex to the Vietoris metric thickening is a homotopy equivalence.
\end{theorem}

The following corollary is immediate by choosing $\cW$ to be as in Section~\ref{ssec:vietoris}.

\begin{corollary}
\label{cor:main}
Let $X$ be a finite-dimensional Polish metric space, and let $r>0$.
The natural bijection $\vr{X}{r}\to\vrm{X}{r}$ is a homotopy equivalence.
For $Z\supseteq X$ a metric space extending the metric on $X$, the natural bijection $\acech{X}{Z}{r}\to\acechm{X}{Z}{r}$ is a homotopy equivalence.
\end{corollary}

The proof of Theorem~\ref{thm:main} will rely on showing that in this setting, $\vietm{\cW}$ is countable-dimensional and an ANR.
The proof that $\vietm{\cW}$ is countable-dimensional is closely related to the fact that if $X$ is a finite-dimensional metric space, then for any integer $k$, the space of all probability measures with support of size at most $k$ is finite-dimensional.
See for example Equation~(3.1) on Page~57 of~\cite{fedorchuk1991probability}, 
%\url{https://www.mathnet.ru/links/7c55e926712f5859b5b382f2bab84f51/rm4568_eng.pdf}
which follows from a more general result of Basmanov~\cite{basmanov1983covariant}.

\begin{lemma}
\label{lem:countable-dimensional}
If $\cW$ is a uniformly bounded open cover of a finite-dimensional Polish metric space $X$, then $\vietm{\cW}$ is countable-dimensional.
\end{lemma}

%(\footnote{\note{Perhaps if $X$ is finite-dimensional, then $\vrm{X}{r}$ is even strongly countable-dimensional, though we don't need that here.}})

\begin{proof}
Recall
\[
\vietm{\cW}=\left\{\sum_{i=0}^{k-1}\lambda_i\delta_{x_i}\in P(X)~\mid~k\ge 1,\ \lambda_i\ge 0,\ \sum_i\lambda_i=1,\ \exists W \in \cW: \{x_0,\ldots,x_{k-1}\} \subseteq W \right\}.
\]
For any integer $k\ge 0$, let
\[
\vietm{\cW}_k=\left\{\sum_{i=0}^{k-1}\lambda_i\delta_{x_i}\in P(X)~\mid~\lambda_i\ge 0,\ \sum_i\lambda_i=1,\ \exists W \in \cW: \{x_0,\ldots,x_{k-1}\} \subseteq W \right\}
\]
be the set of all measures in $\vietm{\cW}$ with support of size at most $k$.
Note that we can write the Vietoris metric thickening as the countable union $\vietm{\cW}=\cup_{k\ge 1}\vietm{\cW}_k$.
So, in order to show that $\vietm{\cW}$ is countable-dimensional, it suffices to show that each $\vietm{\cW}_k$ is finite-dimensional.

Let $Y$ be a topological space.
For $k$ an integer, recall that $P_k(Y)$ is the set of measures with support of size at most $k$,
\[P_k(Y)=\{\mu\in P(Y)~|~\card(\supp(\mu))\le k\},\]
equipped with the weak topology.
%\note{Page~57 of~\cite{fedorchuk1991probability}, 
%%\url{https://www.mathnet.ru/links/7c55e926712f5859b5b382f2bab84f51/rm4568_eng.pdf}
%near Equation~(3.1), states the following.
%``One of the dimensional characteristics of a continuous and intersection preserving functor $\cF\subseteq P_n$ is the formula 
%\[\dim \cF(Y)\le n \dim Y + \dim \cF(n)\]
%that follows from a more general result of Basmanov~\cite{basmanov1983covariant}.
%This formula is valid for any (non-empty) compact space $Y$.''}
%Let $\comp$ be the category of compact metric spaces.
%\note{Question: Is $\cF$ considered as a functor $\cF\colon \comp\to\comp$?
%Do we hence need to add the assumption ``$X$ is compact'' to this lemma and to the subsequent lemma and theorem?}
%\note{Question: Does ``$\cF\subseteq P_n$'' mean $\cF(Y)\subseteq P_n(Y)$ for all $Y\in \comp$?}
%\note{Question: What does ``continuous'' mean here?}
%\note{Question: Does ``intersection preserving'' mean $\cF(\cap_\alpha S_\alpha)=\cap_\alpha \cF(S_\alpha)$?}
For $Y$ a completely regular space, Equation (5.1) of~\cite{fedorchuk1991probability} states $\dim(P_k(Y))\le k \dim(Y)+(k-1)$; see also Basmanov~\cite{basmanov1983covariant}.

Since $X$ is a metric space, it is completely regular, giving $\dim(P_k(X))\le k \dim(X)+(k-1) \eqqcolon D_k$.
Fix a point $a\in X$ and for any positive integer $i\ge 1$, let $X_n = \{x\in X\mid d(a,x)\le i \}$.
Then  the Wasserstein-Kantarovoich-Rubenstein metric on $P_k(X_i)$ agrees with the weak topology on $P_k(X_i)$ by Proposition~\ref{proposition:weaktopology}.
By the subset theorem~\cite[Theorem~5.3.3]{sakai2013geometric}, the inclusion $P_k(X_i)\subseteq P(X)$ gives $\dim(P_k(X_i))\le D_k$ for all integers $i \ge 1$.
Since $P_k(X) =\cup_{i\ge 1} P_k(X_i)$, the countable cum theorem~\cite[Theorem~5.4.1]{sakai2013geometric} implies $\dim P_k(X)\le D_k$.
Finally, the inclusion $\vietm{\cW}_k\subseteq P_k(X)$ gives $\dim(\vietm{\cW}_k)\le D_k$, again by the subset theorem.

%\note{Assuming we are allowed to take $\cF\coloneqq\vr{-}{r}_k\subseteq P_k$ [defined only on \emph{metric spaces}], then we would get
%\begin{align*}
%\dim \vr{X}{r}_k
%&\le k\dim X+\dim\cF(k) \\
%&= k\dim X + (k-1), \\
%\end{align*}
%which would complete the proof.}

Therefore $\vietm{\cW}=\cup_{k\ge 1}\vietm{\cW}_k$ is a countable union of finite-dimensional spaces, and hence is countable-dimensional.
\end{proof}

\begin{lemma}
\label{lem:anr}
If $\cW$ is a uniformly bounded open cover of a finite-dimensional Polish metric space $X$, then $\vietm{\cW}$ is an ANR.
\end{lemma}

\begin{proof}
The space $\vietm{\cW}$ is locally contractible by~\cite[Theorem~2]{HA-FF-ZV}, and it is countable-dimensional by Lemma~\ref{lem:countable-dimensional}.
Hence $\vietm{\cW}$ is an ANR by~\cite[Theorem~6.10.1]{sakai2013geometric}.
\end{proof}

We are now prepared to prove our main result.

\begin{proof}[Proof of Theorem~\ref{thm:main}]
Let $X$ be a finite-dimensional Polish metric space, and let $\cW$ be a uniformly bounded open cover of $X$.
Let $f\colon \viet{\cW}\to\vietm{\cW}$ denote the natural bijection.
We must show that $f$ is a homotopy equivalence

The natural bijection $f\colon\viet{\cW}\to\vietm{\cW}$ is a weak homotopy equivalence by Gillespie's work~\cite{gillespie2024vietoris}.
The space $\viet{\cW}$ is a simplicial complex by definition, and the space $\vietm{\cW}$ is an ANR by Lemma~\ref{lem:anr}.
%Therefore, the combination of Corollaries~6.6.5 and~6.6.6 of~\cite{sakai2013geometric} (as explained in Section~\ref{ssec:Whitehead} on Whitehead's theorem) give that this weak homotopy equivalence is furthermore a homotopy equivalence.
We can now combine Corollaries~6.6.5 and~6.6.6 of~\cite{sakai2013geometric} to get that $f\colon\viet{\cW}\to\vietm{\cW}$ is a homotopy equivalence.
Indeed, let $g\colon Y\to \viet{\cW}$ be a homotopy equivalence from a space $Y$ which is an ANR to the simplicial complex $\viet{\cW}$, which exists by~\cite[Corollary~6.6.5]{sakai2013geometric}.
Then $fg\colon Y\to \vietm{\cW}$ is a weak homotopy equivalence between ANRs, which is therefore a homotopy equivalence by~\cite[Corollary~6.6.6]{sakai2013geometric}.
Let $(fg)^{-1}$ denote a homotopy inverse of $fg$.
Since $g$ and $fg$ are each homotopy equivalences, the two-out-of-three property gives that the natural bijection $f\colon\viet{\cW}\to\vietm{\cW}$ is a homotopy equivalence (with $g(fg)^{-1}$ as a homotopy inverse), as desired.
\end{proof}

Theorem~\ref{thm:main} (or Corollary~\ref{cor:main}) has the following consequences.
The first is Latschev's theorem~\cite{Latschev2001} for metric thickenings.

\begin{corollary}
\label{cor:latschev}
Let $M$ be a closed Riemannian manifold.
There exists $r_0>0$ such that for every $0<r\le r_0$, there is some $\delta>0$ such that for any finite-dimensional Polish metric space $X$ with $d_\gh(X,M)<\delta$, we have $\vrm{X}{r}\simeq M$.
\end{corollary}

\begin{proof}
The result is true with Vietoris--Rips simplicial complexes by~\cite[Theorem~1.1]{Latschev2001}, and then we apply Corollary~\ref{cor:main}.
\end{proof}

The above result resolves \cite[Conjecture~6.9]{AAR} in the affirmative.

\begin{corollary}
\label{cor:interleaving}
Let $X$ be a finite-dimensional Polish metric space.

The natural bijection $\vr{X}{r}\to\vrm{X}{r}$ from the (open) Vietoris--Rips complex to the Vietoris--Rips metric thickening is a $0$-homotopy interleaving, and hence these filtrations have isomorphic persistent homology modules.

For $Z\supseteq X$ a metric space extending the metric on $X$, the natural bijection $\acech{X}{Z}{r}\to\acechm{X}{Z}{r}$ from the (open) \v{C}ech complex to the \v{C}ech metric thickening is a $0$-homotopy interleaving, and hence these filtrations have isomorphic persistent homology modules.
\end{corollary}

We remark that the existence of a $0$-homotopy interleaving~\cite[Definition~3.5]{blumberg2023universality} is a stronger than simply saying that the homotopy interleaving distance is zero --- indeed, it furthermores that the infimum in the definition of the homotopy interleaving distance~\cite[Definition~3.6]{blumberg2023universality} is attained.
When $X$ is a finite-dimensional Polish metric space, Corollary~\ref{cor:interleaving} implies~\cite[Corollary~3]{gillespie2024vietoris}.

\begin{proof}
This is a consequence of the fact that for all $r\le r'$ the following diagrams commute, where the horizontal arrows are inclusions, and where the vertical arrows are the natural bijections, which are homotopy equivalences.
\[
\begin{tikzcd}
	{\vr{X}{r}} & {\vr{X}{r'}} \\
	{\vrm{X}{r}} & {\vrm{X}{r'}}
	\arrow[hook, from=1-1, to=1-2]
	\arrow["\simeq"', from=1-1, to=2-1]
	\arrow["\simeq", from=1-2, to=2-2]
	\arrow[hook, from=2-1, to=2-2]
\end{tikzcd}
\hspace{15mm}
\begin{tikzcd}
	{\acech{X}{Z}{r}} & {\acech{X}{Z}{r'}} \\
	{\acechm{X}{Z}{r}} & {\acechm{X}{Z}{r'}}
	\arrow[hook, from=1-1, to=1-2]
	\arrow["\simeq"', from=1-1, to=2-1]
	\arrow["\simeq", from=1-2, to=2-2]
	\arrow[hook, from=2-1, to=2-2]
\end{tikzcd}
\]
\end{proof}




\section{Strong local contractibility of $\vietm{\cW}$}
\label{sec:strong-local-contractibility}

In this section we prove that for any uniformly bounded open cover $\cW$ of a compact metric space $X$, the space $\vietm{\cW}$ is strongly locally contractible --- in fact strictly strongly locally contractible.
This means, as defined in Subsection~\ref{ssec:local-contractibility}, that each $\mu\in \vietm{\cW}$ has a local base of neighborhoods which strongly deformation retract onto $\mu$.
As before, the cases when $\cW$ is chosen as in Section~\ref{ssec:vietoris} to recover $\vietm{\cW}=\vrm{X}{r}$ or $\vietm{\cW}=\acechm{X}{Z}{r}$ are especially pertinent.

Let $\cW$ be a uniformly bounded open cover of a compact metrizable space $X$.
By $P(X)$ we denote the set of all probability measures on $X$, and by $M(X)$ we denote the space of all measures of $X$.
By the Riesz theorem, $M(X)$ is isomorphic to the dual of $C(X)$, and so $M(X)$ is a Banach space.
In particular, it makes sense to discuss convex subsets of $P(X)$.


\subsection{Alternative basis}

Since $X$ is compact, we have already stated (see Proposition~\ref{proposition:weaktopology}) that the weak topology on $\vietm{\cW}$ coincides with the one induced by the Wasserstein metric.
In this subsection we will construct a convenient basis of this topology, consisting of sets $N(\eta,\cU,\varepsilon)$ defined below.

For a measure $\eta = \sum_{i=1}^n p_i \delta_{y_i}$ on $X$, a continuous functions $f\colon X\to \mathbb R$, and a subset $Y$ of $X$, we let
\[\eta(Y) = \sum_{y_i\in Y} p_i
\quad\mbox{ and }\quad
\eta(f) = \sum_{i=1}^n p_i f(y_i).\]

Recall that the basis of the weak topology on $P(X)$ is given by the sets of the form 
\[\widetilde{O}(\eta, h_1,\dots ,h_k, \varepsilon) =\{\nu \in P(X)~:~ |\eta(h_i)- \nu(h_i)|<\varepsilon \text{ for } i=1,2,\dots, k\},\]
where $h_1, h_2, \dots ,h_k$ are continuous real-valued functions on $X$ and $\varepsilon >0$.
We define
\[O(\eta, h_1,\dots ,h_k, \varepsilon) %= \{\nu \in \vrm{X}{r}~:~ |\eta(h_i)- \nu(h_i)|<\varepsilon \text{ for } i=1,2,\dots, k\}
= \widetilde{O}(\eta, h_1,\dots ,h_k, \varepsilon) \cap \vietm{\cW}.\]

Let $\eta = \sum_{i=1}^n p_i \delta_{y_i}$ be a measure in $\vietm{\cW}$.
Let $\cU=\{U_1, U_2, \dots, U_n\}$ be a collection of disjoint open subsets of $X$ with $y_i\in U_i$ for all $i$, and with $\cup\cU \coloneqq \cup_{i=1}^n U_i$ contained in an element of $\cW$.
For any $\varepsilon >0$, let
\[N(\eta,\cU,\varepsilon) =\left\{\nu \in \vietm{\cW}~:~|\eta(U_i)-\nu(U_i)|<\varepsilon \text{ for all }i\right\}.\]
We will see in Lemma~\ref{lemma:basis3} that the sets of the form $N(\eta,\cU,\varepsilon)$ form a basis for $\vietm{\cW}$ at $\eta$.

\begin{lemma}
\label{lemma:basis1} For each 
\[\nu  =\sum_{j=1}^m w_j \delta_{a_j}\in N(\eta,\cU,\varepsilon)\]
there exists a collection of disjoint open sets $\cV=\{V_1,\dots, V_m\}$ in $X$ and $\delta >0$ such that $\cup\cV$ is contained in an element of $\cW$, such that $a_j\in V_j$ for all $j=1,2,\dots,m$, and such that
\[N(\nu,\cV,\delta)\subseteq N(\eta,\cU,\varepsilon).\]
\end{lemma}

\begin{proof}
Let $\delta = \frac{1}{2m}\min\{\varepsilon - |\eta(U_i) - \nu(U_i)|:i=1,2,\dots,n\}$.
Since $\nu\in \vietm{\cW}$, we have that $\{a_1,\ldots,a_m\}$ is contained in an element of $\cW$.
Consider disjoint open neighborhoods $V_1,\dots , V_m$ of $a_1,\dots , a_m$ such that $V_j\subseteq U_i$ whenever $a_j\in U_i$, and such that $\cup\cV$ is contained in an element of $\cW$.
Let $\cV = \{V_1\ldots,V_m\}$.
Let $\zeta =\sum_{k=1}^l q_k \delta_{c_k}\in  N(\nu,\cV,\delta)$.
Then for all $i=1,2, \dots, n$, we have
\begin{align*}
|\nu(U_i) - \zeta (U_i)|
&\le |\nu(U_i) - \nu(U_i)| + |\nu(U_i) - \zeta(U_i)| \\
&< \varepsilon - 2m \delta + \sum_{\{j\,:\,V_j\subseteq U_i\}} |\nu (V_j) -\zeta (V_j)| + \left|\sum_{c_k\notin \cup\cV} q_k \right| \\
&< \varepsilon - 2m\delta + m\delta + m \delta = \varepsilon.
\end{align*}
Therefore $\zeta \in N(\nu,\cU,\varepsilon)$.
\end{proof}



%Let $\mu =\sum_{i=1}^n m_i \delta_{x_i}$ be a measure in $\vrm{X}{r}$.
%Let $\cU=\{U_1, U_2, \dots, U_n\}$ be a collection of disjoint open subsets of $X$ such that $x_i\in U_i$ for all $i$, and $\diam(\cU) <r$.


\begin{lemma}
\label{lemma:basis2}
For each set of the form $N(\eta,\cU,\varepsilon)$ there exist functions $h_1,h_2,\dots, h_k$ such that
\[O(\eta, h_1,\dots, h_k, \varepsilon)\subseteq N(\eta,\cU,\varepsilon).\]
\end{lemma}

\begin{proof}
Let $\eta = \sum_{i=1}^n p_i \delta_{y_i}$, and let $\cU=\{U_1, U_2, \dots, U_n\}$ be a collection of disjoint open subsets of $X$ with $y_i\in U_i$ for all $i$, and with $\cup\cU$ contained in an element of $\cW$.

Consider functions $f_1,\dots, f_n$ from $X$ to $[0,1]$ such that for all $i=1,2,\dots,n$, we have $f_i(U_i)=\{1\}$ and $f_i(U_j) =\{0\}$ for all $j\ne i$.
Also, consider functions $g_1,\dots, g_n$ from $X$ to $[0,1]$ such that for all $i=1,2,\dots,n$, we have $g_i(y_i) = 1$ and $g_i(X\setminus U_i) =\{0\}$.
We claim that 
\[O(\eta, f_1,\dots, f_n, g_1,\dots, g_n, \varepsilon)\subseteq N(\eta,\cU,\varepsilon).\]
Consider $\nu \in O(\eta, f_1,\dots, f_n, g_1,\dots, g_n, \varepsilon)$.
Note that 
\[\eta(U_i) - \nu (U_i) = \eta(f_i) - \nu (U_i) \le \eta (f_i) - \nu (f_i) <\varepsilon\]
and
\[\eta (U_i) - \nu(U_i) = \eta (g_i) - \nu (U_i) \ge \eta (g_i) - \nu (g_i) > -\varepsilon.\]
This implies $|\eta (U_i) - \nu (U_i)|< \varepsilon$ and therefore $\nu \in N(\eta,\cU,\varepsilon)$.
\end{proof}

\begin{lemma}
\label{lemma:basis3}
The sets of the form $N(\eta,\cU,\varepsilon)$ form a basis for $\vietm{\cW}$ at $\eta$.
\end{lemma}

\begin{proof}
Let $\eta = \sum_{i=1}^n p_i \delta_{y_i}$.
Lemmas~\ref{lemma:basis1} and~\ref{lemma:basis2} imply that sets of the form $N(\nu,\cU,\varepsilon)$ are open in the weak topology.
Consider an open set $O(\nu, h_1,\dots , h_k, \varepsilon)$.
Let
\[M = 1+\max \{ |h_t(x)|~:~ x\in X,\ t=1,2,\dots , k\},\]
and let $\delta = \frac{\varepsilon}{ M(2n+1)}$.
There exists a collection $\cU=\{U_1, \dots , U_n\}$ of disjoint open neighborhoods  of $y_1, \dots, y_n$ such that $\cup\cU$ is contained in an element of $\cW$, and such that for each $i=1,2,\dots, n$ we have $|h_t(y_i) - h_t(x)| < \delta$ for all $x\in U_i$ and $t=1,2,\dots, k$.
Pick $\nu =\sum_{j=1}^m w_j \delta _{a_j} \in N(\eta,\cU,\delta)$.
Then
\begin{align*}
|\eta (h_t) - \nu (h_t)|
&= \left|\sum_{i=1}^n p_i h_t(y_i) - \sum_{j=1}^m w_j h_t(a_j)\right| \\
&\le \left|\sum_{i=1}^n \left( p_i h_t(y_i) - \sum_{a_j\in U_i} w_j h_t(a_j)\right)\right| +\left|\sum_{a_j\notin \cup\cU} w_j h_t(a_j)\right| \\  
&\le \left|\sum_{i=1}^n\left(p_i h_t(y_i) - \sum_{a_j\in U_i} w_jh_t(y_i)\right)\right|+ \sum_{i=1}^n\left(\sum_{a_j\in U_i} w_j |h_t(y_i)-h_t(a_j)|\right)  +M\sum_{a_j\notin \cup\cU} w_j \\
&< M\sum_{i=1}^n \left|p_i - \sum_{a_j\in U_i} w_j\right|+\delta + Mn\delta \\
&\le M\sum_{i=1}^n|\eta(U_i) - \nu (U_i)|+M(n+1)\delta \\
&\le Mn\delta +M(n+1)\delta = \varepsilon.
\end{align*}
 Therefore $\nu\in O(\eta , h_1,\dots, h_k,\varepsilon)$.
\end{proof}

\subsection{Strong local contractibility via Michael's selection theorem for multivalued maps}
\label{ssec:slc}

Our proof that $\vietm{\cW}$ is strictly strongly locally contractible in Theorem~\ref{thm:strong-local-contractibility} will use Michael's selection theorem~\cite{michael1956continuous}, which we recall now.

A multivalued map $F\colon X\to Y$ assigns to each point $x\in X$ a subset $F(x)\subseteq Y$.
A \emph{selection} of $F$ is a function $f\colon X\to Y$ with $f(x)\in F(x)$ for all $x\in X$.

\begin{theorem}[Michael's selection theorem~\cite{michael1956continuous}]
\label{theorem:selection}
Let $X$ be a paracompact space, let $Y$ be a Banach space, and let $F\colon X\to Y$ be a lower semi-continuous multivalued map with non-empty, convex, and closed values.
Then $F$ has a continuous selection $f\colon X\to Y$.
Moreover, if $A$ is a closed subset of $X$ and $g\colon A\to Y$ is a continuous selection of $F|_A$ then $f$ can be assumed to extend $g$.
\end{theorem}

Recall that a multivalued map $F\colon X\to Y$ is \emph{lower semi-continuous} if for any open subset $S$ of $Y$, the set $\{x\in X\mid F(x)\cap S\ne \varnothing\}$ is open in $X$.

We now build up notation necessary for the proof of Theorem~\ref{thm:strong-local-contractibility}.
For any subset $U$ of $X$, we define
\begin{align*}
P_U &= \{ \mu\in \vietm{\cW}\mid \mathrm{supp}\, \mu \cap U\ne \varnothing\} \\
P^U &= \{ \mu\in \vietm{\cW} \mid \mathrm{supp}\, \mu \subseteq U\}.
\end{align*}
Furthermore, for a finite collection $\cU=\{U_1, U_2, \dots, U_n\}$ of subsets $U_i\subseteq X$, let
\[P_{\cU} = \bigcap_{i=1}^n P_{U_i}.\]

Propositions~\ref{prop:open},~\ref{prop:convex}, and~\ref{prop:convex2} have short proofs, which we omit.

\begin{proposition}
\label{prop:open}
If $U$ is an open subset of $X$, then $P_U$ is open in $\vietm{\cW}$.
If $\cU=\{U_1, \dots, U_n\}$ is a collection of open subsets of $X$, then $P_{\cU} $ is open in $\vietm{\cW}$.
\end{proposition}

\begin{proposition}
\label{prop:convex}
For any $Y\subseteq X$, the set $P^Y$ is a convex subset of $P(X)$.
In particular, $P^Y$ is contractible, and every $\mu\in P^Y$ is a strong deformation retract of $P^Y$.
\end{proposition}

Throughout this subsection we fix $\mu =\sum_{i=1}^n m_i \delta_{x_i}$ to be a measure in $\vietm{\cW}$.
Let $\cU=\{U_1, U_2, \dots, U_n\}$ be a collection of disjoint open subsets of $X$ with $x_i\in U_i$ for all $i$, and with $\cup\cU$ contained in an element of $\cW$.
For any $\varepsilon >0$, let
\[ P(\mu,\cU,\varepsilon) = \{ \nu \in \vietm{\cW}~:~\supp(\nu) \subseteq \cup\cU \mbox { and } | \nu(U_i) - \mu(U_i)|\le \varepsilon \mbox{ for all } i \}.\]

\begin{proposition}
\label{prop:convex2}
$P(\mu,\cU,\varepsilon)$ is convex.
\end{proposition}

Let $\tilde{B}(\mu)$ be an open ball in $P(X)$ with respect to the Wasserstein metric, centered at $\mu$, and let $B(\mu) = \tilde{B}(\mu)\cap \vietm{\cW}$.
Note that $\tilde{B}(\mu)$ is convex in $P(X)$, while $B(\mu)$ typically is not: the reason is that, for example, in the Vietoris--Rips metric thickenings, the union of two supports of diameter less than $r$ can have diameter above $r$.
Recall $\cU=\{U_1, U_2, \dots, U_n\}$ is a collection of disjoint open subsets of $X$ such that $x_i\in U_i$ for all $i$, and such that $\cup\cU$ is contained in an element of $\cW$.
We will prove that the open neighborhoods $P_{\cU}\cap B(\mu)$ of $\mu$ are strongly contractible for sufficiently small covers $\cU$.

Note that, for an infinite compactum $X$, the space $P(X)$  can be regarded as a convex compact infinite-dimensional subset of the normed space $\mathbb R^{C(X)}$.
Therefore it is affinely and topologically embeddable in the Hilbert (and hence Banach) space $l_2$ by Klee's theorem~\cite[Theorem~1.1]{klee1955some}.

%\textcolor{red}{Alex: I've added explanation why $P(X)$ can be viewed as a convex subset of $l_2$.}

Define a multivalued map $F\colon P_{\cU}\cap B(\mu) \to P(X)$ by
\begin{equation}
\label{eq:F}
F(\zeta)=\{\nu\in P(\mu,\cU,\varepsilon)~|~\supp(\nu)\subseteq\supp(\zeta)\}.
\end{equation}
 
We will use Michael's selection theorem~\cite{michael1956continuous} to obtain a continuous selection $f\colon P_{\cU}\cap B(\mu) \to P(\mu,\cU,\varepsilon)$.
Thus, we  need to prove that 
the values of $F$ are non-empty, convex, and closed, and that $F$ is lower semi-continuous.

\begin{proposition}
\label{prop:nonempty}
The values of $F$ are non-empty, convex, and closed for any $\zeta\in P_{\cU}\cap B(\mu)$.
\end{proposition}

\begin{proof}
By the assumption, $\supp(\zeta)$  contains a set of the form 
$\{ y^i_1,y^i_2,\dots, y^i_{n_i}\}_{i=1}^n$, where $y^i_j\in U_i$ for each $i$.
Let 
\[\nu =\sum_{i=1}^n \frac{m_i}{n_i} \sum_{j=1}^{n_i} \delta_{y^i_j},\]
where $\mu =\sum_{i=1}^n m_i \delta_{x_i}$.
It is easy to check that $\nu\in F(\zeta)$.

That $F$ is closed-valued follows from the fact that the support of all measures from $F(\zeta)$ is contained in the support of $\zeta$, and from the definition of $P(\mu,\cU,\varepsilon)$.
The convexity follows from Propositions~\ref{prop:convex} and~\ref{prop:convex2}.
\end{proof}

\begin{proposition} \label{prop:lsc}
The multivalued map $F\colon P_{\cU}\cap B(\mu) \to P(X)$ defined in~\eqref{eq:F} is lower semi-continuous.
\end{proposition}

\begin{proof}
To check that $F$ is lower semi-continuous, we must show that for any open subset $S$ of $P(X)$, the set $\{\xi\in P_{\cU}\cap B(\mu)\mid F(\xi)\cap S\ne \varnothing\}$ is open in $P_{\cU}\cap B(\mu)$.
We proceed along the following scheme: for any $\zeta \in P_{\cU}\cap B(\mu)$, any $\nu \in F(\zeta)$, and any arbitrarily small neighborhood $N$ of $\nu$, we define a small neighborhood $N'$ of $\zeta$ such that for each $\xi \in N'$, there exists $\lambda \in F(\xi) \cap N$.

Let
\[\zeta = \sum_{i=1}^n  \sum_{j=1}^{n_i} w^i_j \delta_{y^i_j}+\sum_{t=1}^k w_t \delta_{y_t},\]
where $y^i_j\in U_i$ for all $i,j$, and $y_t\notin \cup\cU$ for all $t=1,2,\dots,k$.
Pick $\nu\in F(\zeta)$.
Next, consider an open neighborhood of $\nu$ which, by Lemma~\ref{lemma:basis3}, may be assumed to be of the form $N=N(\nu,\widetilde \cU,\alpha)$.
Moreover, if 
$\cV=\{V^i_j\mid i=1,\dots,n,\, j=1,\dots, n_i\}$ is a collection of 
%$\{V^i_j\mid i=1,2,\dots ,n, j=1,2,\dots, n_i\}$ are 
disjoint open neighborhoods of the points $y^i_j$ such that $V^i_j\subseteq U_i$ for all $i$, 
and if $\cV'=\{V^i_j\mid y^i_j \in \supp \nu\}$, then
we may further assume that 
$N=N(\nu,\cV',\alpha)$.
%$N$ has the form $N(\nu, \{V^i_j\mid y^i_j \in \supp \nu\},\alpha)$.
Choose $\beta >0$ so that $\beta < w^i_j$ for all $i,j$, and so that
$N' \coloneqq N(\zeta,\cV,\beta)$
%\[N' \coloneqq N(\zeta, \{V^i_j\mid \ i=1,2,\dots,n, j=1,2,\dots ,n_i\},\beta)\]
is contained in $P_{\cU}\cap B(\mu)$, the domain of $F$.
Consider $\xi \in N'$.
Note that the choice of $\beta$ implies that $\supp \xi \cap V^i_j \ne \varnothing$ for all $i,j$.
For each $i=1,2,\dots ,n$ and $j=1,2,\dots, n_i$, pick $z^i_j\in \supp \xi \cap V^i_j$, and let 
\[\lambda = \sum_{
V^i_j\in \cV'
%\nu(V^i_j)\ne 0
} \nu (V^i_j) \delta_{z^i_j}.\]
It is easy to see that $\lambda (U_i) = \nu(U_i)$ for $i=1,2,\dots ,n$, and that $\supp\lambda \subseteq \supp \xi$.
Therefore $\lambda \in F(\xi)$.
Moreover, $\lambda (V^i_j) = \nu (V^i_j)$ for all $V^i_j\in \cV'$, and hence $\lambda \in N$.
Thus, $\lambda \in F(\xi)\cap N$, as required.
\end{proof}

\begin{theorem}
\label{thm:strong-local-contractibility}
If $X$ is a compact metrizable space then $\vietm{\cW}$ is strictly strongly locally contractible.
\end{theorem}

\begin{proof}
As everywhere in this section, let $\mu =\sum_{i=1}^n m_i \delta_{x_i}$ be a measure in $\vietm{\cW}$.
We must show that $\mu$ has a local base of neighborhoods which strongly deformation retract onto $\mu$.

Consider an open ball $B(\mu)$.
Because of Lemma~\ref{lemma:basis3}, we can find $\varepsilon >0$ and a collection $\cU = \{U_1, U_2, \dots, U_n\}$ of disjoint open neighborhoods of $x_1,x_2,\dots, x_n$, such that $\cup\cU$ is contained in an element of $\cW$, and $P(\mu,\cU,\varepsilon)\subseteq B(\mu)$.
By choosing sufficiently small $\varepsilon$ we can also require that $P(\mu,\cU,\varepsilon)\subseteq P_{\cU}$.

Let $F\colon P_{\cU}\cap B(\mu) \to P(X)$ be the multivalued map, defined, as before, by 
\[F(\zeta)=\{\nu\in P(\mu,\cU,\varepsilon)~|~\supp(\nu)\subseteq\supp(\zeta)\}.
\]
By Propositions~\ref{prop:nonempty} and~\ref{prop:lsc}, $F$ is a lower semi-continuous map with nonempty closed and convex values.
Therefore, we can apply Michael's selection theorem~\cite{michael1956continuous} (see Theorem~\ref{theorem:selection}) to obtain a continuous selection $f\colon P_{\cU}\cap B(\mu) \to P(\mu,\cU,\varepsilon)$.
Moreover, by the relative version of the selection theorem, we may assume that $f(\mu) = \mu$.

To construct a deformation retraction of $P_{\cU}\cap B(\mu)$ to $\mu$ that fixes $\mu$, we first join $\zeta$ to $f(\zeta)$ by a straight-line homotopy.
Indeed, note that for all $0\le t\le 1$ we have $(1-t)\zeta+t \cdot f(\zeta)\in \vietm{\cW}\cap \ \tilde{B}(\mu) = B(\mu)$ since $\supp f(\zeta)\subseteq \supp \zeta$, while the continuity of this straight-line homotopy follows from~\cite[Proposition~2.4]{AMMW} or~\cite[Lemma~3.9]{AAF}.
The image of $P_{\cU}\cap B(\mu)$ under this retraction is contained in $P(\mu,\cU,\varepsilon)$, which is convex by Proposition~\ref{prop:convex2}.
And any convex set containing $\mu$ can be strongly deformation retracted onto $\mu$, finishing the proof.
\end{proof}

\begin{corollary}
\label{cor:strong-local-contractibility}
Let $X$ be a compact metric space and let $r>0$.
Then $\vrm{X}{r}$ is strictly strongly locally contractible.
If $Z\supseteq X$ is a metric space extending the metric on $X$, then $\acechm{X}{Z}{r}$ is strictly strongly locally contractible.
\end{corollary}



\section{Conclusion}
\label{sec:conclusion}

We end with some open questions.

\begin{enumerate}

\item For $X$ compact but not necessarily finite-dimensional, is $\vrm{X}{r}$ an ANR or an ANE?

\item Is it true that for $X$ compact, any two sufficiently close maps into $\vrm{X}{r}$ are homotopic?
(This is known to be true for ANRs, and hence is a relaxation of Question (1).)

\item For $X$ an arbitrary metric space, is $\vrm{X}{r}$ strongly locally contractible?

\item When can a map $S^n\to \vrm{X}{r}$ from the $n$-sphere be deformed (up to homotopy) to have its image contained in $\vrm{X}{r} \cap P_{n+1}(X)$?

\item Let $X$ be an ANR.
When is $\vrm{X}{r} \cap P_{n+1}(X)$ an ANR?

%\item Latschev's theorem~\cite{Latschev2001} states that for any $r>0$ sufficiently small, there is a $\delta>0$ such that if $Y$ is $\delta$-close to $M$ in the Gromov--Hausdorff distance, then $\vr{Y}{r'}$ is homotopy equivalent to $M$ for all $r'<r$.
%This is also true for the metric thickening $\vrm{Y}{r'}$ when $Y$ is finite, simply because for $Y$ finite the Vietoris--Rips simplicial complex and metric thickening are homeomorphic.
%Is there a version of Latschev's theorem with metric thickenings that allows for infinite $Y$?
%\note{We can now answer affirmitively, using Corollary~\ref{cor:main}, when $Y$ is Polish and finite-dimensional!
%Could still keep a question about a setting with fewer assumptions.}

\end{enumerate}



\bibliographystyle{plain}
\bibliography{VietorisThickeningsAndComplexesOfManifoldsAreHomotopyEquivalent}



\end{document}

